% !TEX root = ../linnik399.tex
\section{Concluding remarks}\label{sec:remarks}

\emph{The difficult configurations.} The tightest certified cases have a nonreal first
character with $\lambda_1\in[0.64,0.80]$ and a nearby second family. The numerical
experiments of \cref{subsec:why} suggest that stronger location bounds in this range
would be useful. One limitation of a single near row can be seen directly: if $N$
entries have the same positive feature $v$ and diagonal $D$, then its quadratic
inequality permits $N\le D/(v^2-d)$ when $v^2>d$ (\cref{ex:twolevels}). This is a restriction of that
relaxation, not a lower bound on what other tests or arguments can achieve.

Improved far-density constants may also help. Xylouris estimates
\cite[p.~87]{Xylouris2011nullstellen} that halving the constant in his (5.19) would lower
the exponent obtained by his method to about $4.6$. A corresponding gain for the
present method would require new certificates for the complete cover.

\emph{Possible improvements.} Heath-Brown's list of possible improvements
\cite[\S16]{HeathBrown1992zero} contains several that we have not used: optimal test functions
for the combined inequalities (his item~1), positivity beyond the support of the test functions via
Brun--Titchmarsh bounds (item~3; see \cite{Titchmarsh1930divisor,MontgomeryVaughan1973large,
Iwaniec1982brun,Maynard2013brun}), and modified sieve coefficients in the density argument (item~6;
compare \cite{FriedlanderIwaniec2023selberg}). His item~7, a continuous weight in the density
argument, is what Xylouris's far-density lemma implements, and we use it. For cube-free moduli Burgess's bounds hold for every $k$, and even the
Weyl bound is known \cite{PetrowYoung2020weyl}; Heath-Brown already noted that all the arguments
improve in that case. Recent large-value estimates for Dirichlet polynomials
\cite{GuthMaynard2026new,ChenGuptaLi2025large,TaoTrudgianYang2026new} and new density theorems
\cite{Pintz2019some} improve zero-density estimates away from $\sigma=1$; it is not clear
whether they help in the range $\sigma=1-O(1/\log q)$ that matters here.

\emph{Other approaches.} For moduli with bounded cubic part the exponent $4.5$ has been claimed
\cite{Meng2010note}, and for moduli all of whose prime factors are small much better
exponents are known \cite{Chang2014short}. A different recent approach, through the exceptional set of Goldbach's problem (compare
\cite{Pintz2018new}), claims the exponent $5$ in a preprint \cite{Zhao2025exceptional}.

\emph{The role of the computation.} The graded lemma gives quadratic constraints for
different test functions and anchors. Their threshold forms can be combined with the
far-density budget in linear programs, whose bounds are certified by exact duality.
This separates the search for effective test parameters from the verification of the
resulting bound. Any improvement of the global exponent must still include the
zero-location arguments, a complete case cover, and the exterior regimes.
